\section{Closing}

% ---------- Slide 18 ----------
\begin{frame}{They also built a simulation tool}
  \SectionPill{Tool}{PresFourAccent}
  \begin{columns}[T,onlytextwidth]
    \begin{column}{0.50\textwidth}
      \centering
      \begin{tikzpicture}[node distance=0.25cm]
        \node[modelbox,text width=1.8cm,minimum height=0.85cm] (a) {input\\graph};
        \node[modelbox,text width=1.8cm,minimum height=0.85cm,right=0.25cm of a] (b) {configure\\moves};
        \node[modelbox,text width=1.8cm,minimum height=0.85cm,right=0.25cm of b,
              fill=PresFourAccentSoft,draw=PresFourAccent] (c) {run\\simulations};
        \node[modelbox,text width=2.4cm,minimum height=0.85cm,below=0.55cm of b] (d) {save \& visualize};
        \draw[-Latex,line width=0.6pt] (a)--(b);
        \draw[-Latex,line width=0.6pt] (b)--(c);
        \draw[-Latex,line width=0.6pt] (c.south) -- ++(0,-0.27) -| (d.east);
      \end{tikzpicture}
      \vspace{0.55em}

      {\scriptsize\color{PresFourMuted}Built on Python + NetworkX.}
    \end{column}
    \begin{column}{0.47\textwidth}
      \begin{beamercolorbox}[wd=\linewidth,rounded=true,sep=0.50em]{takeaway box}
        {\color{PresFourAccent}\bfseries What the tool does}\par
        \small
        \begin{itemize}\setlength\itemsep{0.10em}
          \item generate synthetic networks
          \item import real-world networks
          \item configure perturbation criteria
          \item search possible graph changes
          \item save and visualize results
        \end{itemize}
      \end{beamercolorbox}
      \vspace{0.30em}
      {\small\color{PresFourMuted}Framed as a reusable framework rather than a one-off experiment.}
    \end{column}
  \end{columns}
  \SlideNote{1:30}{Reproducibility. The tool matters because the paper frames this as a framework others can apply.}
\end{frame}

% ---------- Slide 19 ----------
\begin{frame}{What should we be careful about?}
  \SectionPill{Limitations}{PresFourAdversary}
  \begin{columns}[T,onlytextwidth]
    \begin{column}{0.32\textwidth}
      \MiniCard{PresFourAdversary}{1. Scale}{Mostly small synthetic networks (20--100 nodes).\\ Exhaustive search is hard on large graphs.}
    \end{column}
    \begin{column}{0.32\textwidth}
      \MiniCard{PresFourAdversary}{2. Realism}{Some global moves may be socially expensive or impossible.\\ Fake nodes and edge deletions have very different real-world costs.}
    \end{column}
    \begin{column}{0.32\textwidth}
      \MiniCard{PresFourAdversary}{3. Model scope}{Single adversarial node.\\ Mostly undirected, unweighted, static graphs.\\ Real platforms are temporal and directed.}
    \end{column}
  \end{columns}
  \vfill
  \begin{beamercolorbox}[wd=\linewidth,rounded=true,sep=0.55em]{takeaway box}
    {\color{PresFourAccent}\bfseries Open question}\quad
    Would the same strategies work on real social media data, or would platform constraints absorb them?
  \end{beamercolorbox}
  \SlideNote{2:30}{Show that you understand the contribution but can evaluate it critically.}
\end{frame}

% ---------- Slide 20 ----------
\begin{frame}{Influence metrics are attack surfaces}
  \SectionPill{Takeaway}{PresFourAdversary}
  \centering
  \vspace{0.10cm}
  {\fontsize{16}{20}\selectfont\bfseries Small graph edits can cause large changes in perceived power.}\par

  \vspace{0.55cm}
  \begin{columns}[T,onlytextwidth]
    \begin{column}{0.32\textwidth}
      \MiniCard{PresFourAccent}{1.}{Observed networks are fragile.}
    \end{column}
    \begin{column}{0.32\textwidth}
      \MiniCard{PresFourAdversary}{2.}{Centrality rankings can be gamed.}
    \end{column}
    \begin{column}{0.32\textwidth}
      \MiniCard{PresFourDecrease}{3.}{Robustness should be part of SNA.}
    \end{column}
  \end{columns}

  \vspace{0.35cm}
  \begin{columns}[c,onlytextwidth]
    \begin{column}{0.78\textwidth}
      \begin{beamercolorbox}[wd=\linewidth,rounded=true,sep=0.55em]{takeaway box}
        {\color{PresFourAdversary}\bfseries Discussion}\par
        \small If centrality can be flipped by a single move, should platforms or researchers report it as a \emph{stable} measure of influence?
      \end{beamercolorbox}
    \end{column}
    \begin{column}{0.20\textwidth}
      \centering
      \begin{tikzpicture}[scale=0.9]
        \node[gadv] (a) at (0,0) {};
        \node[gnode] (n) at (1.2,0.0) {};
        \draw[addedge,-{Latex[length=2mm]}] (a) to[bend left=20] (n);
      \end{tikzpicture}
    \end{column}
  \end{columns}
  \SlideNote{1:30}{End with an informal discussion prompt rather than only summarizing.}
\end{frame}
